\section{The active three-plane and ambient amplification}
All cross products in this section are taken solely in the oriented Euclidean
three-plane $E$. They define an orthonormal basis of $E$; they are never used
as a divergence or pressure identity on $\R^d$. Write $M=\nabla u(t,0)$ for
the parent central gradient and $H=\nabla^2P(t,0)$. Oddness fixes the central
particle at zero. Reflection makes $M,H$ block diagonal on $E\oplus E^\perp$.
The active block need not be trace free.

The geometric input at a stage is the decomposition
\begin{equation}\label{eq:central-decomposition}
 M=B+hqp^T+\mathcal E,\qquad |\mathcal E|\leq E_*,\qquad h(t_0)=h_*,
\end{equation}
where $p,q$ are the normalized transported normal and primary velocity of
the preceding packet. They lie in $E$ and are orthogonal; put $n=p\times q$.
The shear is the actual preceding packet amplitude,
$h=c_0|m_{\rm old}||v_{\rm old}|$ with constant $c_0>0$.
Consequently $\dot h/h=-B_{pp}-B_{qq}$, by the two transport equations
for the preceding normal and velocity.
Here $B$ is the gradient of the preceding background solution.
On the stage interval $|B|\leq C G_*$ and $|\dot B|\leq C G_*^2$.
The transported normal and velocity satisfy
\begin{equation}\label{eq:ray-velocity}
 \dot m=-M^Tm,\qquad
 \dot v=-Mv+2m\frac{m\cdot Mv}{|m|^2},\qquad m\cdot v=0.
\end{equation}
The coefficient 2 follows by differentiating the last constraint. These
are full ambient equations; for central data in $E$ they preserve $E$.

At activation put
\begin{gather*}
 a=p\cdot Bq\in[1/2,2],\quad \beta=(n\cdot Bq)/a,
 \quad \tfrac12\leq\beta x_{\rm prev}^2\leq2,\\
 \epsilon=\sqrt{a/h_*},\quad
 \tau=a(t-t_0)/\epsilon,\quad x=\sqrt\beta\,\tau.
\end{gather*}
The target is $x=x_{\rm tar}\geq2$. Let
$\Theta\geq2+\beta^{-1/2}+x_{\rm tar}/\sqrt\beta$ bound the scaled time,
with at most $\Theta^{-60}$ scaled time retained after the target.
The geometric error parameter is
\begin{equation}\label{eq:geometric-error}
 e=C_d\bigl(\epsilon\Theta G_*^2+E_*+
                   \rho K_h^{c_d}/\epsilon\bigr),
 \qquad e\Theta^{60}\leq c_d x_{\rm prev}^{-10}.
\end{equation}
Here $\rho$ is the physical packet length and $K_h$ bounds the actual
parent particle map, its inverse deformation and required derivatives.
The small constant on the right is fixed before the scales are chosen.

\begin{lemma}[Endpoint selection]\label{lem:actual-endpoint}
On a positive history interval $[0,t_0]$, assume the one-sided
pressure bound $H\leq K_+I$ and $K_+t_0^2/2\leq1/2$, provided by
the mean reserve. Suppose $\|M\|_\infty\leq C_Mh_*$,
$\|H\|_\infty\leq C_Hh_*h_{\rm old}$, $1\leq h_{\rm old}\leq h_*$,
and $t_0^{-1}\leq h_*$. Suppose also
$\overline B=B(t_0)+\mathcal E(t_0)$ has
$\overline B_{pp}<0$ and $\|\overline B\|\leq\zeta h_*$.
For sufficiently small $\zeta$, depending only on $C_M,C_H$, an endpoint
$Y\in E\cap n^\perp$ of size at most $C/h_*$ gives a stationary displacement
whose terminal primary velocity satisfies
\[
 v_q=1,\qquad -C\leq v_p\leq0,\qquad v_{E^\perp}=0.
\]
The displacement is stationary in the full transverse space.
\end{lemma}
\begin{proof}
Choose $m_0=F(t_0,0)^Tn/|F(t_0,0)^Tn|$ and transport it by $F^{-T}$.
Let $\Lambda$ be the endpoint operator of Lemma~\ref{lem:endpoint-symmetry}.
Integration by parts in its stationary equation gives
\[
 Y\cdot\Lambda Z=\int_0^{t_0}
    (\eta_Y'\cdot\eta_Z'-H\eta_Y\cdot\eta_Z)\,dt.
\]
The stationary paths start at zero. The inequality
$\int|\eta|^2\leq(t_0^2/2)\int|\eta'|^2$ therefore makes their
action nonnegative even for a nonzero terminal value. Thus $\Lambda$
is symmetric and positive semidefinite. To bound it, put
$L_0=t_0^{-1}+\|M\|_\infty+\|H\|_\infty^{1/2}+1$.
A comparison path vanishes before $t_0-L_0^{-1}$ and then equals
$L_0(t-t_0+L_0^{-1})\operatorname{proj}_{m(t)^\perp}Y$.
Since $\|\partial_t\operatorname{proj}_{m^\perp}\|\leq C\|M\|$,
its action is at most
$C(L_0+\|M\|_\infty^2/L_0+\|H\|_\infty/L_0)|Y|^2
\leq C_0h_*|Y|^2$.
Minimality yields $0\leq\Lambda\leq C_0h_*I$ on the whole transverse
space, and reflection reduces its restriction to $E\cap n^\perp$ to a
symmetric two-by-two matrix.

Write $b=\Lambda_{pq}$, $d_0=\Lambda_{qq}$ and
$u_*=\Lambda_{pp}-\overline B_{pp}>0$. The terminal velocity
$v=\eta'-M\eta$ has active coordinates
\[
 \binom{v_p}{v_q}=
 \begin{pmatrix}u_*&b-\overline B_{pq}\\
 b-h_*-\overline B_{qp}&d_0-\overline B_{qq}\end{pmatrix}
 \binom{Y_p}{Y_q}.
\]
If $b\leq h_*/2$, take $Y=-p$. Then $v_p=-u_*\leq0$ and
$v_q=h_*+\overline B_{qp}-b\geq h_*/3$.
If $b>h_*/2$, positivity implies
$\Lambda_{pp}\geq b^2/d_0\geq h_*/(4C_0)$.
Take $Y=q-(b-\overline B_{pq})p/u_*$. Its norm is bounded by a
constant depending on $C_0$, and $v_p=0$. Moreover
\[
 v_q=\left(d_0-\frac{b^2}{u_*}\right)-\overline B_{qq}
 +\frac{bh_*+b(\overline B_{pq}+\overline B_{qp})
       -h_*\overline B_{pq}-\overline B_{qp}\overline B_{pq}}{u_*}.
\]
The parenthesis is nonnegative because $u_*\geq\Lambda_{pp}$.
The numerator is at least $h_*^2/2-C(C_0+1)\zeta h_*^2$.
Since $u_*\leq(C_0+\zeta)h_*$, choosing $\zeta$ small gives
$v_q\geq c h_*$. Divide the displacement, endpoint and velocity by
this positive terminal value. Both cases give the asserted normalization
and endpoint bound. Reflection gives exactly zero extra components
before and after normalization. At $t_0=0$ the prescription is simply
$v=q$, and no stationary history problem is needed.
\end{proof}

\begin{lemma}[The scalar growing solution]\label{lem:scalar-growth}
Let $0<\beta\ll1$, $P_0=\beta\tau^2$, $Q_0=-2\beta\tau$, and
$D=1+P_0^2$. The equation
\[
 (DV')'=2(1-\beta P_0)V,\qquad V_\lambda(0)=1,
 \qquad V_\lambda'(0)=\lambda\geq0
\]
has a positive solution to every target under consideration, and
$V_\lambda(x=1)\geq e^{b/\sqrt\beta}$. For $x\geq1$ put
$r_\lambda=-V_\lambda'/V_\lambda$. There is a bounded nonnegative $z$
such that
\[
 r_\lambda=\frac{\sqrt\beta}{x}-\frac z{x^2},\qquad
 z^2=\frac2{1+x^{-4}}+O(\sqrt\beta)\quad(x\geq2).
\]
Writing $g(\tau)=V_0(\tau)$, the two-component propagator for
$(-V',V)$ obeys
\[ \|\Phi_0(t,s)\|\leq C\Theta^8g(t)/g(s). \]
All statements concern a scalar equation and are dimension independent.
\end{lemma}
\begin{proof}
On $x\leq1$, integrating twice with $D\leq2$ and
$2(1-\beta P_0)\geq1$ gives
$V_\lambda\geq1+\lambda\tau/2+\tfrac12\int_0^\tau(\tau-s)V_\lambda(s)ds$.
Positive iteration bounds it below by $\cosh(\tau/\sqrt2)$.
Its logarithmic derivative $l$ satisfies $0\leq l$ and
$l'\leq2-l^2$, whence $l\leq\sqrt2\coth(\sqrt2\tau)$.
For $x\geq1$, first write $\widehat V(x)=V(x/\sqrt\beta)$,
and then put $s=1-\rho$, $\rho=x^{-1}$ and
$f(\rho)=\widehat V(1/\rho)/\rho$. Direct substitution gives
\[
 ((1+\rho^4)f_\rho)_\rho=(2/\beta-2\rho^2)f.
\]
Starting from $f(1)>0$, $f_\rho(1)<0$, backward integration preserves
both signs. For $z=-\sqrt\beta f_\rho/f$ and
$\mu=\sqrt{(2-2\beta\rho^2)/(1+\rho^4)}$,
\[
 \sqrt\beta z_s=\mu^2-z^2+
               \sqrt\beta\frac{4\rho^3}{1+\rho^4}z.
\]
This gives a uniform upper bound and preserves nonnegativity.
Indeed its initial value at $\rho=1$ is
$z(1)=\sqrt\beta+V_\lambda'(1/\sqrt\beta)/V_\lambda(1/\sqrt\beta)$,
which is bounded independently of $\lambda$ by the preceding
logarithmic-derivative estimate.
The equation for $z-\mu$ has damping coefficient $z+\mu\geq c>0$
and a bounded forcing multiplied by $\sqrt\beta$. Its integrating
factor gives $|z-\mu|\leq Ce^{-c(1-\rho)/\sqrt\beta}+C\sqrt\beta$,
which proves the direction formula for $\rho\leq1/2$.

Reduction of order gives
$V_\lambda=g(1+\lambda I)$, where $I(\tau)=\int_0^\tau(Dg^2)^{-1}$
and $I(1)\geq c>0$. Positivity before $x=1$, and increase of $xg$
afterwards, give $g(s)/g(u)\leq C\Theta$ for $s\leq u$.
For arbitrary scalar data at $s$,
\[
 V(t)=\frac{g(t)}{g(s)}\left[
 V(s)+D(s)\left(V'(s)-\frac{g'(s)}{g(s)}V(s)\right)
 \int_s^t\frac{(g(s)/g(u))^2}{D(u)}\,du\right].
\]
The integral is at most $C\Theta^3$, $D(s)\leq C\Theta^4$, and
$|g'/g|\leq C$. Near $\tau=0$ the latter bound follows from
$g'(0)=0$ and $l'\leq2-l^2$; away from zero it follows from the
displayed bounds for $l$ and $z$.
Differentiating this formula costs at most one further
power of $\Theta$, proving the propagator assertion. This is the scalar
calculation in \cite[\S\S4.4--4.5]{OAI}, reproduced to specify exactly the
part of that argument used here.
\end{proof}

The selected initial conditions at activation are
$m(t_0)=s_0n(t_0)$, $v_q(t_0)=1$, $v_n(t_0)=0$ and
$-C\leq v_p(t_0)\leq0$. At the central label express
$m=s_0(Pp+\epsilon Qq+Nn)$ and
$v=\epsilon Up+Vq+\epsilon Wn$. The moving-frame connection has
$(S_f)_{12}=B_{pq}$, $(S_f)_{13}=B_{pn}$ and
$(S_f)_{32}=-B_{nq}$; skew symmetry determines its other entries.
Equations~\eqref{eq:ray-velocity} then give
\[
 (P,Q,N)=(\beta\tau^2,-2\beta\tau,1)+O(e\Theta^5),
 \qquad W=-(PU+QV)/N,\quad N\geq1/2.
\]
The reference ray system is $P'=-Q$, $Q'=-2\beta N$, $N'=0$.
Its triangular propagator costs $C\Theta^2$. The exact error integral
is at most $Ce\Theta^5+Ce\Theta^3\sup|\mathrm{error}|$, and the last
term is absorbed by \eqref{eq:geometric-error}. Substitution in the
velocity equation gives a coefficient error at most $Ce\Theta^{13}$
from the two-component ideal system
\[
 U'=-2V+\frac{2P_0((P_0+\beta)V+Q_0U)}{1+P_0^2},\qquad V'=-U.
\]
Indeed the only unsuppressed contributions to
$J=(m\cdot Mv)/(s_0a)$ are $PB_{pq}V/a$, $h\epsilon^2QU/a$
and $NB_{nq}V/a$. The identity $\dot h/h=-B_{pp}-B_{qq}$ gives
$h/h_*=1+O(\epsilon\Theta G_*)$ on this interval, so
$h\epsilon^2/a=1+O(\epsilon\Theta G_*)$. The remaining contributions
contain either $\epsilon G_*$ or the error in
\eqref{eq:geometric-error}. Transversality bounds
$|W|\leq C\Theta^2(|U|+|V|)$; the denominator $|m|^2/s_0^2$
is at least $1/4$. These observations bound every rational coefficient
without using a trace condition on the active block.
The full matrix estimate, including all transverse coordinates, is
derived in Section~\ref{sec:parameter-envelope}.

Weighting Duhamel by $g(s)/g(t)$, Lemma~\ref{lem:scalar-growth} makes
its error operator at most $Ce\Theta^{22}$. A geometric series therefore
gives the exact active propagator bound $2C\Theta^8g(t)/g(s)$.
For the selected data, $\lambda=-v_p(t_0)/\epsilon\geq0$,
\[
 |U-U_\lambda|+|V-V_\lambda|
       \leq Ce\Theta^{30}(1+\lambda)g.
\]
For $\tau\geq1$, $V_\lambda\geq c(1+\lambda)g$, so this is a
relative error. In particular
\begin{equation}\label{eq:positive-flux}
 J/V=P_0+\beta+Q_0r_\lambda+O(e\Theta^{34})>0.
\end{equation}
For $x\leq1$ its leading term is at least $\beta$, since
$Q_0,r_\lambda\leq0$. For $x\geq1$ it equals
$x^2-\beta+2\sqrt\beta z/x$. The error is less than half the
lower bound by \eqref{eq:geometric-error}.

\begin{lemma}[Central frame renewal]\label{lem:central-renewal}
At the target set $p_2=m/|m|$, $q_2=v/|v|$, $n_2=p_2\times q_2$,
$a_2=p_2\cdot Mq_2$ and $\beta_2=(n_2\cdot Mq_2)/a_2$. Then
\begin{align*}
 a_2/a&=1+O(x_{\rm tar}^{-4}+\beta x_{\rm tar}^{-2}
       +\sqrt\beta x_{\rm tar}^{-3}+e\Theta^{40}),\\
 \beta_2x_{\rm tar}^2&=1+O(\sqrt\beta+x_{\rm tar}^{-4}+e\Theta^{40}),\\
 p_2\cdot Mp_2&\leq-\frac{c\sqrt{h_*}}{x_{\rm prev}x_{\rm tar}}
                                     +C(G_*+E_*).
\end{align*}
Moreover $A_{\rm tar}=|m||v|$ at the target satisfies
$A_{\rm tar}\geq K_h^{-c}\Theta^{-c}e^{b x_{\rm prev}}$.
\end{lemma}
\begin{proof}
Put $r=U/V$, $w=W/V$, $D_a=|m|^2/s_0^2$ and
$E_v=1+\epsilon^2(r^2+w^2)$. Normalize the two vectors directly:
\[
 a_2/a=\frac{J/V}{\sqrt{D_a}\sqrt{E_v}},\quad
 \beta_2=\frac{[(P,\epsilon Q,N)\times(\epsilon r,1,\epsilon w)]
                    \cdot(M(v/V)/a)}{(J/V)\sqrt{E_v}}.
\]
In the second numerator substitute
$P=x^2$, $Q=-2\sqrt\beta x$, $r=\sqrt\beta/x-z/x^2$.
The terms of size $\beta x^2$ and $\sqrt\beta xz$ cancel, leaving
\[
 -1+(1+x^{-4})z^2+\beta x^{-2}-2\sqrt\beta z x^{-3}
       =1+O(\sqrt\beta).
\]
The perturbation contributes at most $Ce\Theta^{36}$.
Also $J/V=x^2-\beta+2\sqrt\beta z/x+O(e\Theta^{34})$,
$D_a=1+x^4+O(e\Theta^7)$ and $E_v=1+O(\epsilon^2\Theta^4)$.
These identities prove the first two estimates. The rank-one part of
\eqref{eq:central-decomposition} gives
$p_2\cdot Mp_2=h\epsilon QP/D_a+O(G_*+E_*)$.
Since $QP\leq-c\sqrt\beta x^3$, $D_a\leq Cx^4$ and
$h\epsilon\asymp\sqrt{h_*}$, this proves compression.
Finally $V_\lambda(x_{\rm tar}/\sqrt\beta)\geq x_{\rm tar}^{-1}
V_\lambda(1/\sqrt\beta)\geq x_{\rm tar}^{-1}e^{b/\sqrt\beta}$ and
$s_0^{-1}\leq K_h^c$; the relative-error estimates give the lower bound
for $A_{\rm tar}$.
\end{proof}

Let $\mathcal P\geq1$ bound the source quantities
$K_h,\Theta,\epsilon^{-1},h_{\rm child},\delta^{-1}$ and the fixed
coefficient and inverse costs. These quantities are estimated in
Section~\ref{sec:parameter-envelope} before the frequency is selected.
The choice $\alpha=\delta h_{\rm child}/A_{\rm tar}$ therefore satisfies
$\alpha\leq\mathcal P^c e^{-b x_{\rm prev}}$.
At the central target the cutoff equals one, the phase is zero and
$f_\delta'(0)=\delta^{-1}$, so the leading gradient is exactly
$\alpha\delta^{-1}v\otimes m=h_{\rm child}q_2p_2^T$.
Section~\ref{sec:parameter-envelope} proves the corresponding full
transverse propagator and pressure-sign estimates in a neighborhood of
the central label. The following consequence of that comparison gives
a bound uniform over the stages of the iteration.

\begin{lemma}[Uniform target-size comparison]\label{lem:uniform-packet-size}
For the selected primary, on every label in the quantitative packet
neighborhood, the full ambient relative estimates of
Section~\ref{sec:parameter-envelope} give
\[
 \frac{|m(t,y)||v(t,y)|}{A_{\rm tar}}\leq C_d
       \qquad(1\leq\tau\leq\tau_{\rm end}).
\]
On the history interval and on $0\leq\tau<1$ the same ratio is at most
$K_h^{c_d}\Theta^{c_d}\epsilon^{-c_d}e^{-b x_{\rm prev}}$.
For $g(\tau)=V_0(\tau)$ and
$\alpha=\delta h_{\rm child}/A_{\rm tar}$, one also has
$\sup\alpha g\leq\mathcal P^{c_d}$ and
$\alpha\leq\mathcal P^{c_d}e^{-b x_{\rm prev}}$.
\end{lemma}
\begin{proof}
The full relative comparison, including its extra components, and
$\epsilon\Theta^2\ll1$ imply, for $\tau\geq1$,
\[
 |m(t,y)||v(t,y)|\asymp_d
       s_0\sqrt{1+x^4}\,V_\lambda(x/\sqrt\beta).
\]
The small perturbations of the initial label and of $s_0$ are included
in the same constants. At the central target this equivalence bounds
$A_{\rm tar}$ from below by a fixed multiple of the right side.
These constants do not grow with $\Theta$ or the stage.

For $0\leq x\leq1$, $V_\lambda$ is increasing. For $x\geq1$ the
reciprocal-variable proof of Lemma~\ref{lem:scalar-growth} shows that
$xV_\lambda(x/\sqrt\beta)$ is increasing. Hence for
$1\leq x\leq x_{\rm tar}$,
\[
 (1+x^2)V_\lambda(x/\sqrt\beta)
 \leq (x+x^{-1})x_{\rm tar}
             V_\lambda(x_{\rm tar}/\sqrt\beta)
 \leq2x_{\rm tar}^2V_\lambda(x_{\rm tar}/\sqrt\beta).
\]
For $x\leq1$ the same conclusion follows by comparison first with
$V_\lambda(1/\sqrt\beta)$ and then with the target. Since
$\sqrt{1+x^4}\asymp1+x^2$, this proves the constant size comparison
up to the target. Afterwards the logarithmic derivative in $\tau$ of
$\sqrt{1+x^4}V_\lambda$ is
\[
 \frac{2\sqrt\beta x^3}{1+x^4}-\frac{\sqrt\beta}{x}
                                          +\frac z{x^2},
\]
which is bounded for $x\geq x_{\rm tar}\geq2$. The retained scaled
interval has length at most $\Theta^{-60}$, so its extra factor is at
most two. The full relative errors preserve this fixed bound.

The stationary history inverse bounds the unnormalized history by
$K_h^{c_d}$. On $\tau\leq1$, the scaled system has bounded
coefficients and initial norm at most $C(1+\lambda)$, with
$\lambda\leq C/\epsilon$; converting the coordinates therefore costs
only the displayed old-data polynomial. Divide these estimates by
the exponential target lower bound from
Lemma~\ref{lem:central-renewal} to obtain the early and history ratio.
Finally $V_0\leq V_\lambda$. The monotonicity argument just used gives
$\sup g\leq C x_{\rm tar}V_\lambda(x_{\rm tar}/\sqrt\beta)$,
whereas
$A_{\rm tar}\geq c_d s_0x_{\rm tar}^2
V_\lambda(x_{\rm tar}/\sqrt\beta)$.
Their growing factors cancel in $\alpha g$. The remaining factors
$\delta,h_{\rm child},s_0^{-1},x_{\rm tar}$ have the fixed polynomial
source bounds. The target exponential lower bound similarly proves
the assertion for $\alpha$ itself.
\end{proof}
